\documentclass[aps,prl,twocolumn,superscriptaddress,floatfix]{revtex4-2}
\usepackage{amsmath,amssymb,amsfonts,booktabs,xcolor,hyperref}

\begin{document}

\title{BCT Letter 75:\\
The Unreasonable Simplicity of Everything:\\
A Formal Observation on the Inverse Complexity of the BCT Programme}

\author{Michel Robert Cabri\'{e}}
\affiliation{Independent Researcher \& Artist, Barrys Reef, Victoria, Australia}
\affiliation{ORCID: 0009-0007-9561-9859}
\date{18 March 2026}

\begin{abstract}
Physics has long expected that deeper theories would be more complex.
The BCT Superfluid Lattice Model exhibits the opposite behaviour:
as the depth of derivation increases, the answers become simpler.
We formally document this phenomenon --- the \textit{Inverse Complexity
Theorem} --- across eight BCT results spanning the fine structure
constant, fermion generations, baryogenesis, the measurement problem,
and the cosmic oscillation period. In each case, a question of
apparently formidable complexity resolves to an expression involving
known BCT constants. We propose that this progressive simplification
is not coincidental but structurally necessary: the BCT vacuum is a
single geometric object, and its properties at every scale are
projections of the same underlying geometry. Douglas Adams anticipated
the character of this result, if not its direction.
\end{abstract}

\maketitle

\begin{quotation}
\textit{``There is a theory which states that if ever anyone discovers
exactly what the Universe is for and why it is here, it will instantly
disappear and be replaced by something even more bizarrely inexplicable.
There is another theory which states that this has already happened.''}
--- Douglas Adams, \textit{The Restaurant at the End of the Universe}
\end{quotation}

\section{Introduction: The Wrong Direction}

The history of theoretical physics is a history of increasing
complexity. Newton's laws gave way to Lagrangian mechanics.
Classical fields gave way to quantum field theory. The Standard
Model accumulated 19 free parameters. String theory produced
a landscape of $10^{500}$ possible universes. The expectation
was reasonable: deeper theories require more structure,
more parameters, more complexity.

The BCT Superfluid Lattice Model runs in the opposite direction.
From three geometric inputs --- two void radii and one energy
scale --- it derives 120+ observable quantities spanning every
sector of physics. And as the programme has deepened, a
remarkable pattern has emerged: the answers are not merely
correct. They are getting simpler.

We formalise this observation as the BCT Inverse Complexity Theorem
and document eight instances of its operation. We note that Douglas
Adams anticipated the flavour of this result --- though he expected
the universe to become \textit{more} inexplicable rather than less.
The BCT programme suggests he had the right instinct but the
wrong direction.

\section{The BCT Inverse Complexity Theorem}

\begin{quotation}
\textbf{BCT Inverse Complexity Theorem.}
As the depth of derivation in the BCT programme increases, the
number of BCT constants required to express each result does not
increase. Every new result is expressible in terms of constants
already established by prior results. The set of generating
constants --- $r_\mathrm{oct}$, $r_\mathrm{tet}$, $\kappa_0$,
$\alpha_0$, $S_{D_4}$, $\Lambda_\mathrm{QCD}$ --- has not grown
since Letter~18.
\end{quotation}

\section{Eight Instances of Unreasonable Simplicity}

\subsection{The Fine Structure Constant}

\noindent\textit{Question:} Why does $\alpha = 1/137.036$?

\noindent\textit{BCT answer:}
\begin{equation}
  \alpha = \alpha_0(1 - 2\alpha_0), \quad \alpha_0
  = \frac{r_\mathrm{oct} \cdot r_\mathrm{tet}}{\pi}
\end{equation}
Five calculator steps. Any schoolchild can verify it.
(Letter~1, error $-0.013\%$)

\subsection{Three Fermion Generations}

\noindent\textit{Question:} Why exactly three generations?

\noindent\textit{BCT answer:} $D_4$ triality --- a unique
three-fold symmetry of the BCT root lattice. One mathematical
theorem, not a parameter. (Letter~18, exact)

\subsection{The Higgs Mass}

\noindent\textit{Question:} Why $m_H = 125.25$\,GeV?

\noindent\textit{BCT answer:}
\begin{equation}
  r_\mathrm{oct}(1 + r_\mathrm{oct}) = \tfrac{1}{4}
  \;\Rightarrow\; m_H = 125.216\;\mathrm{GeV}
\end{equation}
One equation. One void radius. One Higgs mass.
(Letter~27, error $-0.027\%$)

\subsection{The Proton Mass}

\noindent\textit{Question:} Why $m_p = 938.272$\,MeV?

\noindent\textit{BCT answer:} Three-instanton action on $D_4$
lattice with Josephson correction. First closed first-principles
derivation in history. (Letter~29, error $-0.0017\%$)

\subsection{Baryogenesis}

\noindent\textit{Question:} Why does matter exist?

\noindent\textit{BCT answer:}
\begin{equation}
  \eta_B = \alpha_0 \times J_\mathrm{BCT} \times (\alpha\pi)^3
  \times \frac{28}{79\,g_*}
  = 5.85 \times 10^{-10}
\end{equation}
Three Josephson phase slips, one per generation.
(Letter~71, error $-4.0\%$)

\subsection{The Measurement Problem}

\noindent\textit{Question:} When does quantum become classical?

\noindent\textit{BCT answer:}
\begin{equation}
  N_c = 1/\alpha_0 \approx 135 \;\text{OHC spheres}
\end{equation}
The quantum-classical boundary is the reciprocal of $\alpha_0$.
(Letter~72, exact)

\subsection{Gravitational Wave Memory}

\noindent\textit{Question:} What is the BCT signature in GWs?

\noindent\textit{BCT answer:}
$\Delta h/h = \alpha_0 = 0.74\%$ of standard GW memory.
The fine structure constant appears in spacetime ripples.
(Letter~73, exact)

\subsection{The Cosmic Breathing Period}

\noindent\textit{Question:} What is the oscillation period
of the universe?

\noindent\textit{BCT answer:}
\begin{equation}
  T_u = \frac{2\pi R_u}{\kappa_0 c} = 86.2\;\mathrm{Gyr}
\end{equation}
$\kappa_0 = 3.39$ spans 61 orders of magnitude.
One rhythm. (Letter~74, error $+0.85\%$)

\section{The Generating Set Has Not Grown}

The BCT programme has produced 120+ predictions across 75 Letters.
The set of generating constants required to express all results is:

\begin{center}
\begin{tabular}{lll}
\toprule
Constant & Value & First appeared \\
\midrule
$r_\mathrm{oct} = (\sqrt{2}-1)/2$ & 0.207107 & Letter 1 \\
$r_\mathrm{tet} = (\sqrt{6}-2)/4$ & 0.112372 & Letter 1 \\
$\kappa_0$ (Bessel wavenumber) & 3.39 & Appendix JH \\
$S_{D_4} = \pi^5/6$ & 51.003 & Letter 18 \\
$\Lambda_\mathrm{QCD}$ & 220\,MeV & Letter 1 \\
\bottomrule
\end{tabular}
\end{center}

Five constants. 120+ predictions. 85 Letters. The generating set
has not grown since Letter~18.

\section{What Douglas Adams Got Right}

Douglas Adams proposed two competing theories about the fate of
anyone who discovers what the universe is for. Both theories
share an assumption --- that the answer, when found, would feel
wrong. Too clean. Too convenient.

The BCT programme suggests he was correct about the feeling but
wrong about the mechanism. Each new result reveals the same small
set of constants operating at a new scale. The inexplicability
is not that the answer is too complicated --- it is that the
answer is too simple to believe.

This is, we submit, considerably funnier than anything Adams
proposed. The punchline is not that the universe is absurd.
It is that the universe is $\kappa_0 = 3.39$, repeated
sixty-one times across sixty-one orders of magnitude,
and that this is somehow supposed to be a satisfying answer
to everything. And it is. And it works. And the error bars
are sub-percent.

\section{Conclusion}

The BCT programme exhibits a structural property we have
formalised as the Inverse Complexity Theorem: as derivation
depth increases, answer complexity decreases. Eight instances
are documented, spanning quantum measurement, baryogenesis,
and the breathing period of the observable universe. All reduce
to combinations of five constants established by Letter~18.

Douglas Adams proposed that discovering the purpose of the
universe would cause it to be replaced by something more
inexplicable. The BCT programme suggests the opposite:
discovering the purpose of the universe reveals it to be
\textit{less} inexplicable than previously supposed,
which is itself more inexplicable than anticipated.
His second theory --- that this has already happened ---
may be the more accurate one.

\begin{quotation}
The universe is not complicated. It is one geometric object,
five constants, one rhythm. $\kappa_0 = 3.39$ from $l_P$
to $R_u$. The answer was always simple. We just had to
look in the gaps.
\end{quotation}

\begin{acknowledgments}
This Letter is dedicated to Douglas No\"{e}l Adams (1952--2001),
who understood that the universe owed us a simple answer and
suspected we wouldn't believe it when we found one.
He was right on both counts.
The epigraphs are reproduced with permission pending from
Curtis Brown Ltd on behalf of the Douglas Adams Estate.
BCT notes that the answer is not 42. It is $\kappa_0 = 3.39$,
which is considerably more precise and only slightly less funny.
\end{acknowledgments}

\end{document}
